% Section 1: Introduction
% Note: \section{} header is in main.tex

Iterative self-training with reward-based curation has become a prominent strategy for language model alignment. The current model generates multiple candidate responses per prompt, a reward model scores them, and the top-scoring response is selected for supervised fine-tuning to produce the next-generation model. This loop raises a fundamental stability question: if each generation retains only a narrow fraction of outputs (1 out of 4 or 1 out of 8 candidates), the compounding effect should progressively narrow the output distribution and erode the breadth that makes language models practical across diverse user needs.

Two research threads support this concern. The model collapse literature demonstrates that training on synthetic outputs leads to progressive distribution narrowing~\cite{placeholder_shumailov, placeholder_alemohammad}: rare modes are underrepresented in synthetic samples, this compounds across iterations, and the model collapses onto high-density modes. Separately, RLHF studies document that single-round reward optimization reduces lexical and semantic variety~\cite{placeholder_kirk}. Together, these establish a strong prior: iterative reward-curated self-training should cause diversity collapse with increasing severity at each generation.

Yet neither literature directly addresses the deployed setting. Model collapse results study unfiltered self-consumption where the model trains on its full output distribution; reward-curated pipelines apply an external selection filter, retain a fixed prompt scaffold across generations, and decode stochastically, all of which violate the theoretical preconditions for progressive narrowing. The RLHF literature examines single rounds rather than multi-generation iteration. No prior work has systematically measured whether diversity degrades across multiple generations of reward-curated fine-tuning under controlled factorial conditions.

We fill this gap with a controlled experiment spanning 8 configurations and 3 generations. The central result overturns expectations: semantic diversity (Vendi Score~\cite{placeholder_vendi}) remains stable across all curated conditions ($-0.27\%$ to $+0.91\%$), including extreme selection pressure ($K{=}16$, 6.25\% retention, $\Delta = -0.19\%$). However, this stability masks \emph{metric-dependent collapse}: format diversity degrades substantially regardless of curation, with near-duplicate start rate rising from 37.9\% to 75.6\% and opening entropy dropping 0.75 bits at $K{=}4$. Standard embedding-based metrics are blind to this convergence. A second finding concerns the interaction between adaptation method and curation: under LoRA, curation suppresses Vendi Score by $-3.40$ pp relative to random, but under Full-FT with higher learning rate, curation promotes diversity by $+0.63$ pp. This pattern is consistent with a capacity artifact: LoRA's rank-16 bottleneck limits how much distributional structure the model can absorb, and the apparent suppression disappears when this constraint is removed.

Our experimental design isolates four factors: adaptation method (LoRA versus Full-FT), selection pressure ($K \in \{4, 8, 16\}$, corresponding to 25\%, 12.5\%, and 6.25\% retention rates), learning rate ($5 \times 10^{-6}$ versus $2 \times 10^{-5}$), and curation strategy (reward top-1 versus random selection). All experiments use Qwen3-0.6B (596M parameters) as the base model, ArmoRM-Llama3-8B-v0.1~\cite{placeholder_armorm} as the reward model, and a fixed set of 2{,}200 prompts from UltraFeedback. This factorial structure enables direct attribution of observed diversity changes to specific experimental factors rather than confounding variation in data composition or model scale.

Our contributions are as follows:
\begin{enumerate}
    \item We provide systematic evidence that reward curation does not cause semantic diversity collapse in iterative self-training. Across 8 configurations and 3 generations, all curated conditions maintain Vendi Score within 1\% of the initial value, even under extreme selection pressure ($K{=}16$, 6.25\% retention). This result contradicts expectations from the model collapse literature.
    \item We identify metric-dependent collapse: format diversity (response openings, structural templates) degrades substantially while embedding-based semantic diversity remains stable. This degradation occurs at comparable rates under curated and random selection, which indicates it is inherent to iterative self-training rather than caused by reward filtering. The effect intensifies with candidate pool size ($K{=}16$, $-$0.19\% Vendi but 1.13 bits opening entropy loss at $K{=}8$). We propose two complementary indicators, \texttt{opening\_entropy} and \texttt{near\_dup\_start\_rate}, that capture this previously invisible dimension of homogenization, and recommend their inclusion in diversity audits alongside standard embedding-based measures.
    \item We observe a direction reversal in the effect of curation across adaptation methods: curation suppresses diversity under LoRA ($-3.40$ pp) but promotes it under Full-FT ($+0.63$ pp), a pattern consistent with LoRA's rank-16 bottleneck constraining how much distributional structure the model can absorb from curated data.
\end{enumerate}